% Folder 78, batch 02 — archivists' pages 21–40.
% Pass: Opus 5.5 (claude-opus-5-5), 2026-10-02 — first pass, unchecked against the pages by a human.
%
% Notes for maintainers: E is the affine plane carrying the polygon, Aff(E)
% its affine group; D_n is the dihedral group of order 2n, D_n^0 its rotation
% subgroup, u its generating rotation, s_0 the base vertex, Pi_0(n) the
% standard combinatorial n-gon, Pol(S) the set of polygon structures on S,
% Psi_n = (Z/nZ)^*/{+-1}. Sections IV to VII of his numbering fall in this
% batch; the formula numbers (14)-(28) are his. Page 21 opens mid-sentence and
% page 40 breaks off mid-sentence. Writing is fast throughout: formulas carry,
% connective prose often does not.
\documentclass[11pt,a4paper]{article}
\input{../preamble/grothendieck}

\folder{78}
\batch{2}
\pages{21}{40}
\foldertitle{[Polygones réguliers et polynômes cyclotomiques] : notes manuscrites (s.d.).}
\dating{[à partir de 1977]}
\watermark{Édition de démonstration}

\begin{document}

\page{21}
\note{la page s'ouvre au milieu d'une phrase commencée avant ce lot ; les renvois (11), (12), (13), (19) et aux propositions 1 à 3 visent des pages antérieures.}

donne une bijection \add{((19) disait que c'était une surjection)} entre l'ens.\ des isom.\ $\mathrm{Isom}((E,\varphi),(E',\varphi'))$ et l'ens.\ des $n$-pol.\ \uncertain{rég.} \add{\uncertain{ép.}} sous-jacents à~: $\varphi'$.

D'une façon suggestive, on peut dire que la donnée d'un polygone $n$-régulier sous-jacent à une repr.\ $\varphi$ admissible (ce qui revient au même, le choix d'un pt $s \in D_n^{*}$) revient à la donnée d'un \uncertain{épinglage} de $\varphi$ (ayant comme effet d'éliminer les automorphismes --- ou $G$ opère librement sur $D_n^{*}$ --- sans introduire de nouvelles données d'isom.\ --- ou $G$ opère transitivement ---).
\note{l'indice de $D^{*}$ est mal formé sur cette page ; on le lit $n$ d'après la p.~24, où $D_n^{*} \setminus D_n^{0}$ est écrit nettement.}

\page{22}
\section{IV. Où on laisse tomber les conditions de fidélité}

Considérons des données (11) satisfaisant (13) mais pas néc.\ (12)~; donc~: pas l'unimodularité~; on laisse tomber la condition $\varphi$ fidèle. \uncertain{Comment} se présente-t-elle~?

\textbf{Proposition 4} \struck{Pour que $\varphi$ soit fidèle, il faut et il suffit que}

(a) Si $v \in D_n^{0}$, i.e.\ $v = u^{i}$, alors
\[
  v_E = \mathrm{id} \iff v(s_0) = s_0 .
\]

(b) Pour que $\varphi$ soit fidèle, il faut et il suffit que $\varphi \mid D_n^{0}$ soit fidèle, \struck{\ill{}} (ou encore) que
\[
  u^{i} s_0 = s_0 \implies i \equiv 0 \ (n) \qquad (\text{i.e. } u^{i} = 1),
\]
ou encore que $\forall d \in \mathbb{N}^{*}$, $d \mid n$, on ait $u^{d}(s_0) \neq s_0$.
\note{sous $\mathbb{N}^{*}$ un signe isolé, peut-être une retouche (« $d \neq n$ » ?), non lu.}

(c) \struck{fait} Le noyau de $\varphi$ \struck{\ill{}} est engendré par un (unique) élément $u^{d}$, où \struck{\ill{}} $d$ est un diviseur de $n$ tel que \add{$n' = n/d \geq 3$} \struck{$d \geq 3$.} \struck{\ill{} tel que $\varphi$ \ill{}} \struck{est induit \ill{} par \ill{}}

\textit{Dém.} (a)
\[
  v_E = \mathrm{id} \iff v(s_\nu) = s_\nu \quad \forall \nu
\]
(car les $s_\nu$ engendrent affinement $E$) or $s_\nu = u^{\nu} s_0$, donc $v(s_\nu) = v u^{\nu} s_0 = u^{\nu} v s_0$, donc il suffit $v(s_0) = s_0$.

\page{23}
\uncertain{(c)} Le noyau de $\varphi \mid D_n^{0}$ est un sous-groupe de $D_n^{0} = \mathbb{Z}/n\mathbb{Z}$, donc \struck{de la} c'est le s-groupe engendré par un $u^{i}$, où $i$ est un diviseur $> 0$ déterminé de $n$ --- et que a priori pourrait être $1$ ou $2$. Mais si \struck{$i$} $n'$ est tel que
\[
  n = n' d
\]
\struck{$\varphi$ se factorise en} \struck{$\varphi' : D_{n'} \to \mathrm{Aff}(E)$} on aura que \ill{} $\mathrm{card}\, S = n'$, et comme $S$ engendre aff.\ $E$, on aura $n' \geq 3$. Enfin, comme $\varphi$ se factorise par
\[
  \varphi' : D_{n'} \longrightarrow \mathrm{Aff}(E)
\]
avec $\varphi' \mid D_{n'}^{0}$ fidèle --- et que les conditions (11) sont évidemment satisfaites aussi par $\varphi'$ --- (b) implique que $\varphi'$ est fidèle, donc $\mathrm{Ker}\, \varphi$ est le même que $\mathrm{Ker}(\varphi \mid D_n^{0})$.

(b) \struck{Résulte de $\varphi \mid D_n^{0}$ équivalent à \ill{} fidèle.}
\marginal{remarque sur (c)}
\note{l'étiquette du premier alinéa est mal formée ; la démonstration est celle de (c), et elle s'appuie sur (b), qui reste sans démonstration propre.}

\page{24}
\textbf{Corollaire 1} Si $v \in D_n^{*} \setminus D_n^{0}$, alors $\varphi(v) \neq \mathrm{id}_E$.

En effet, \struck{car} il est clair que $\varphi(v)$ échange les deux ordres circulaires \add{$\{v_E \mid v \in \omega\}$} \struck{sur} $S$, \struck{distincts de l'opposé \ill{}}, donc ne peut être l'identité sur $S$, \struck{donc $v_E$} donc pas non plus sur $E$.

\textbf{Corollaire 2} Considérons la situation \add{du (b)} ci-dessus. Alors \struck{$\varphi' : D_{n'}$}
\[
  \begin{cases}
    \varphi' : D_{n'} \longrightarrow \mathrm{Aff}(E) \\
    s_0 \in E
  \end{cases}
\]
définissent dans $E$ un polygone épinglé à $n'$ côtés.

\textbf{Scholie} On peut considérer que la donnée $\{\varphi, s_0\}$ \struck{correspond}~: une réalisation géom.\ d'un polygone comb.\ type $\Pi_0(n)$, qui est un revêtement à $d$ feuillets de la réalisation géom.\ \add{$(\varphi', s_0)$} \struck{des} du polyg.\ comb.\ type $\Pi_0(n')$.

Quand on fixe des \uncertain{cohérents} avec un

\page{25}
$n$ fixé, on trouve \add{d'abord} des \struck{\ill{}} repr.\ $\varphi : D_n \to \mathrm{Aff}(E)$, sans savoir d'avance quand elles sont fidèles --- ceci serait une question d'un autre type que les précédentes \add{\ill{}} (trouver les \uncertain{solutions} satisfaisant les conditions (13)), demandant une analyse plus \struck{\ill{}} délicate de la situation. Surtout quand on regardera des réalisations géom.\ d'un $\Pi_0(n)$ sur un anneau plus général qu'un corps, tel que $\mathbb{Z}$ ou mieux, les anneaux tels que $\mathbb{Z}[T]/(T^{n}-1)$, ou les anneaux cyclotomiques --- on trouvera \add{pour les fibres \ill{}} \struck{\ill{}} les \ill{} \ill{} \ill{} \ill{} (y compris en car.\ \ill{}) une repr.\ fidèle, tandis qu'en \ill{} caractéristiques elles devraient être fidèles (\uncertain{ou fidèles ?})) --- elles n'en cessent pas

\page{26}
pour autant d'exister, et de définir des polygones réguliers en un \ill{} \ill{}, mais à un nb plus petit de côtés (diviseur de $n$). Ce serait emprisonnant --- et \add{\uncertain{finalement d'ailleurs}} artificiel --- de traîner tout au long dans les données la condition de fidélité, au lieu de se borner à l'examiner seulement aux moments \struck{voulus} où cette considération est utile.

\section{V. \ill{} des $\infty$-polygones}
\note{le premier mot du titre est récrit sur un mot biffé ; ni l'un ni l'autre n'est lu.}

On peut aller encore plus loin, en notant que la donnée de (11) (a) revient~: la donnée de
\[
  (14)\qquad
  \begin{cases}
    \sigma_E,\ \sigma'_E \in \mathrm{Aff}(E) \quad \text{satisfaisant} \\
    \sigma_E^{2} = \sigma_E'^{2} = \mathrm{id}_E ,
    \quad (\underbrace{\sigma'_E \sigma_E}_{u_E})^{n} = \mathrm{id}_E
  \end{cases}
\]
\struck{satisfaisant} \struck{(15) la fidélité signifiant que (\ill{})} et la donnée de $s_0$ \struck{la fidélité} satisfaisant \struck{(15)} (13)
\[
  \begin{cases}
    \sigma_E s_0 = s_0 \\
    s_0,\ s_1 = u_E s_0 \ \text{et}\ s_2 = u_E^{2} s_0 \ \text{sont non alignés}
  \end{cases}
\]
\note{dans la seconde ligne, avant $s_1 = u_E s_0$, une première écriture biffée ($\sigma_E \ldots$) n'est pas reproduite.}

\page{27}
la condition de fidélité s'écrit
\[
  (15)\qquad \forall \text{ diviseur } d \mid n, \ \text{on a} \ u_E^{d}(s_0) \neq s_0
\]
\add{$u_E^{d} \neq \mathrm{id}$ i.e.} (cf prop 4 (b)) --- et il suffit de l'exiger pour $3 \leq d \leq \left[\frac{n}{3}\right]$ (donc si $3 \leq n \leq 8$, cette condition est automatiquement satisfaite, ce premier entier où elle ne l'est pas est $n = 9$, où il suffit de \uncertain{tester} pour $n = 9$ --- le $9$-gone régulier \struck{\ill{}} pouvant dégénérer en \uncertain{trois fois} un $3$-gone régulier ---). --- Mais on peut fort bien laisser tomber cette condition de fidélité, et laisser tomber \add{\uncertain{même}} la condition (14)
\[
  (\sigma'_E \sigma_E)^{n} = \mathrm{id}_E
\]
ce qui signifie qu'on regarde en fait une représentation
\[
  (16)\qquad \varphi : D_\infty \longrightarrow \mathrm{Aff}(E)
\]
du groupe diédral infini, sans se préoccuper si elle se factorise par $D_n$ --- étant entendu qu'

\page{28}
on s'intéressera plus particulièrement, parmi les actions en question, aux repr.\ de $D_\infty$ qui \uncertain{veulent} bien se factoriser par $D_n$. C'est dire que, \struck{\ill{}} de ce point de vue, la question de fidélité ne se pose plus --- les représentations \add{\uncertain{peut-être}} les plus intéressantes de $D_\infty$ \struck{\ill{}} --- celles qui donnent naissance aux polygones réguliers finis --- étant précisément non fidèles~!

\marginal{Pt de vue $\Pi$-épinglé (\ill{} $\infty$-polygones \ill{} \ill{})}
\note{note marginale écrite en oblique, reliée par un trait à « fidèles » ; lecture très incertaine.}

Arrivé à ce point, il y a lieu de considérer --- de même que des polygones \add{$\Pi = (S, \omega)$} combinatoires infinis (tous connexes \add{ou repr.}) polyg.\ comb.\ finis --- des réalisations géom.\ de ceux-ci (qui sont fidèles), et de même des repr.\ géométriques régulières, \struck{\ill{}} \struck{\ill{} \ill{} \ill{}} qui sont \uncertain{données} par la donnée d'une

\page{29}
réalisation géom.\ satisfaisant les conditions~:
\[
  (17)\qquad \forall v \in \mathrm{Aut}(\Pi), \quad \exists !\, v_E \in \mathrm{Aut}(E)
\]
telle que $v_E(\varphi(s)) = \varphi(v(s))$

\begin{tikzcd}
  S \arrow[r, "\varphi"] \arrow[d, "v"'] & E \arrow[d, "v_E"] \\
  S \arrow[r, "\varphi"'] & E
\end{tikzcd}

\struck{Ce que} La donnée d'une telle réalisation revient~: la donnée
\[
  (18)\qquad
  \begin{cases}
    \text{a)}\ \ \varphi_\Gamma : \Gamma = \mathrm{Aut}(\Pi) \longrightarrow \mathrm{Aff}(E),
      \quad v \longmapsto v_E \quad \text{homom.\ de groupes} \\
    \text{b)}\ \ S \xrightarrow{\ \varphi_S\ } E
  \end{cases}
\]
\add{(application compatible avec les opérations de $\Gamma$)} satisfaisant les \add{seules} conditions supplémentaires
\[
  (19)\qquad S_E \ \text{n'est pas contenu dans une droite}
\]
\struck{(20) \ill{}}
\note{à gauche de (19), un griffonnage biffé (une première numérotation (19), (20)) n'est pas reproduit.}

Si on se fixe d'avance $i_0 \in S$, et si $\sigma_0$ désigne l'unique élt.\ de $\Gamma$ $\neq 1$ tel que \struck{$\sigma_0 i_0 = i_0$}
\[
  (20)\qquad \sigma_0 \neq 1_\Gamma , \quad \sigma_0 i_0 = i_0
\]
la donnée de $\varphi_S$ équivaut \add{(pour $\varphi_\Gamma$ donné)} à la donnée de
\[
  (21)\qquad
  \begin{cases}
    s_0 \in E & [\, s_0 = \varphi_S(i_0) \,] \\
    \text{satisfaisant} & \sigma_0 s_0 = s_0
  \end{cases}
\]
et la condition (19) équivaut, donc, \struck{\ill{}} si on choisit $u \in \omega$, à
\[
  (22)\qquad \text{\struck{$s_0 \neq s_1$, $s_0 \in E$}} \quad
  s_0,\ s_1 = u s_0,\ s_2 = u^{2} s_0 \quad \text{non alignés.}
\]

\page{30}
Mais la donnée de $(s_0, u)$ équivaut~: celle d'un \emph{épinglage} de $\Pi$, i.e.\ d'un \struck{autre $\Pi$ de} isom.
\[
  \Pi \simeq \Pi_0(n)
\]
(où $n$ \struck{\ill{}} est un entier $\geq 3$ ou $+\infty$), et la donnée de a) \struck{(a)} équivaut donc~: la donnée de (14) (sauf la dernière \add{condition de (14), si $n = +\infty$}) \struck{condition}, et celle de (b) équivaut~: la donnée de $s_0 \in E$ satisfaisant la première relation $\sigma_E s_0 = s_0$, et la condition (19) se traduisant par (22). D'où résumé~:

\textbf{Prop 5} Les \struck{repr.} réalisations géom.\ \add{régulières} \add{dans $E$} de $\Pi_0(\infty)$ correspondent aux \struck{couples} triples $(\sigma_E, \sigma'_E, s_0)$ de deux automorphismes affines \struck{de} de $E$
\[
  \sigma_E,\ \sigma'_E \in \mathrm{Aff}(E)
\]
et d'un \struck{$s_0 \in E$, satisfaisant les}
\[
  s_0 \in E
\]
satisfaisant les conditions
\begin{align*}
  &(1^{\circ}) && \sigma_E^{2} = \sigma_E'^{2} = \mathrm{id} \\
  &(2^{\circ}) && \sigma_E s_0 = s_0 \\
  &(3^{\circ}) && s_0,\ s_1 = u_E s_0,\ s_2 = u_E^{2} s_0 \ \text{non alignés}
\end{align*}
(où l'on a posé $u_E = \sigma'_E \sigma_E$).
\note{en (3°), devant $s_0$, une formule surchargée et biffée n'est pas reproduite.}

\page{31}
\struck{et enfin la condition d'unimodularité} \struck{$(4^{\circ})$ $u_E$ unimodulaire}

On dit que la réalisation régulière est \emph{unimodulaire} si elle satisfait
\[
  (4^{\circ})\qquad u_E \ \text{est unimodulaire.}
\]

\textbf{Scholie} En étendant ces objets, on étudie simultanément tous les $n$-polygones réguliers \add{épinglés} pour $n \in \mathbb{N}^{*}$, $n \geq 3$ --- ils correspondent aux cas où $u_E$ est d'ordre fini, et $n$ est l'ordre de $u_E$ --- et même davantage qu'un nouveau type d'objets~: les $\infty$-polygones réguliers épinglés.

\textbf{Remarque} Les considérations finales concernant les polygones réguliers s'étendent aux $\infty$-polygones réguliers sans \ill{} dégâts. Notamment \uncertain{tous} prop.~1, et corollaire, la prop.~2 (qui s'étend en l'énoncé évident \struck{avec}) en ce qui concerne la dernière question de fidélité $D_\infty \hookrightarrow \mathrm{Aff}(E)$ \uncertain{correspond} \struck{conditions} \uncertain{finies}. Bien sûr $\infty$-pol.\ rég.

\marginal{Mais il conviendrait de \uncertain{donner} la définition des polygones réguliers généraux (\uncertain{finis ou} \struck{\ill{}} \add{\ill{}}) au lieu de la définition d'abord donnée, à cas finis …}
\note{note marginale écrite en oblique dans la marge gauche du bas de la page, raccordée par un trait à « prop. 1 ».}

\page{32}
la prop.~3 \struck{\ill{}} montrant le lien exact entre $\infty$-polygones réguliers et repr.\ affines du groupe diédral infini, enfin la prop.~4 (où dans (b) \add{et (c)}, il ne faut plus parler de divisibilité, puisque $n$ est ici remplacé par $\infty$ …)

\textbf{NB} J'ai \uncertain{omis}, pour une donnée d'\uncertain{iso.}\ des réal.\ géom.\ \add{régulières} \struck{des} polygones combinatoires, \struck{\ill{}} \struck{il y a \ill{}} « \ill{} » $\alpha \in k$, \uncertain{maintenant} $\alpha \in k$ n'est plus soumis à aucune condition, cyclotomique ou autre --- ) il y a lieu de \uncertain{dire} « la » repr.\ aff.\ sur $E$ correspondant \emph{au} représentant épinglé (ant.\ unique). Plus général et plus intrinsèque, si $\Pi$ est un polyg.\ combinatoire fini ou infini, notons que la cat.\ de ses réalisations géom.\ régulières est \emph{rigide} (sans automorphismes autres que l'identité) --- il y a lieu par là de parler d'une (\uncertain{ou de la} $= \alpha$) réalisation \add{$= E(\Pi)$}, $E$ --- étant entendu que \uncertain{si on sait} une

\marginal{C'est un des \uncertain{intérêts techniques} d'éliminer \struck{\ill{}} au départ \uncertain{de la théorie} …}

\page{33}
\uncertain{réalisation} intrinsèque, de façon que l'opération de $\Gamma \simeq \mathrm{Aut}(\Pi)$ sur $E(\Pi)$ sera obtenue automatiquement. (Je \uncertain{montre} \uncertain{donc simplement}~: décrire $S \longrightarrow E(\Pi)$ compatible avec l'op.\ de $\Gamma$ …).

\section{VI. Opérations $\Psi_r$ d'étoilage d'un polygone}

Soit $(S, \omega)$ un polygone combinatoire, et soit $r \in \mathbb{Z}$, considérons
\[
  (23)\qquad \Psi_r(\omega) = \{\, u^{r} \mid u \in \omega \,\}
\]
\struck{A' que}

\textbf{Prop 6} a) Pour que \struck{\ill{}} $(S, \Psi_r(\omega))$ soit un polygone combinatoire (i.e.\ pour que $u^{r}$ opère transitivement sur $S$) il faut et il suffit que \struck{$S$} $\mathrm{card}(S) = n$ soit fini, et que $(r, n) = 1$. \struck{Si $\mathrm{card}(S) = n$,}

b) Si \struck{$\mathrm{card}\, S = n$} $n$ est fini, \struck{\ill{}} \struck{\ill{} $\Psi_r(\omega)$} alors pour $r, r' \in \mathbb{Z}$, on a
\[
  \Psi_r(\omega) = \Psi_{r'}(\omega) \iff r' \equiv \pm r \ (n)
\]
\struck{\ill{} \ill{} $\mathbb{Z}^{*}/\pm \longrightarrow \mathrm{Pol}(S)$} \add{donc $\Psi_n$ opère librement} sur $\mathrm{Pol}(S)$.

c) On a, pour $r, r' \in \mathbb{Z}$
\[
  \Psi_r(\Psi_{r'}(\omega)) = \Psi_{rr'}(\omega)
\]
i.e.\ $(\mathbb{Z}/n\mathbb{Z})^{*}$ \uncertain{opère} sur l'ens.\ $\mathrm{Pol}(S)$ des structures de polygones généraux sur $S$. En fait, c'est
\[
  (24)\qquad (\mathbb{Z}/n\mathbb{Z})^{*}/\{\pm 1\} \overset{\text{déf}}{=} \Psi_n
\]
\uncertain{pour} $\mathrm{Pol}(S)$, qui opère \uncertain{librement}.
\note{la formule (24) et la fin de c) sont écrites dans la marge gauche et en tête de la marge, reliées par une flèche ; le texte biffé autour de b) est très surchargé et seule sa partie lisible est rendue.}

\page{34}
\textit{Dém.} \struck{\ill{}} Soit $u \in \omega$. La condition que $(S, \omega)$ soit un polyg.\ comb.\ est que $\mathbb{Z}$ opère (via $u$) transitivement sur $S$ \add{et $n \geq 3$}, i.e.\ revient à que $S$ soit un torseur sous $\mathbb{Z}_n$ (où $\mathbb{Z}_n = \mathbb{Z}/n\mathbb{Z}$ si $n$ fini, $\mathbb{Z}_\infty = \mathbb{Z}$). \struck{Dans l'opération de} L'opération de $\mathbb{Z}$ sur $S$ déduite des $u^{r} \in \Psi_r(\omega)$ n'est autre que celle qui se factorise par $\mathbb{Z}_n$ et se déduit de l'op.\ précédente via l'homom.
\[
  \mathbb{Z}_n \xrightarrow{\ r\,\mathrm{id}_{\mathbb{Z}_n}\ } \mathbb{Z}_n
\]
Donc elle fait de $S$ un $\mathbb{Z}_n$-torseur ssi $r\,\mathrm{id}_{\mathbb{Z}_n}$ est un iso, ce qui signifie que $n$ fini et $(r, n) = 1$.

(b) \struck{\ill{} $\Psi_r(\omega) = \Psi_{r'}(\omega)$}

\struck{\ill{}} (c) $\Psi_r(\omega) = \Psi_{r'}(\omega)$ signifie
\[
  u^{r'} = u^{r} \quad \text{ou} \quad u^{r'} = (u^{-1})^{r} \ (= u^{-r})
\]
donc
\[
  r' \equiv r \ (n) \quad \text{ou} \quad r' \equiv -r \ (n).
\]
\note{l'alinéa marqué (c) démontre en fait b) de la Prop.~6.}

\textbf{Remarque} Une structure de $n$-pol.\ \emph{orienté} sur $S$ équivaut~: la donnée d'une permutation circulaire $u$ sur $S$,

\page{35}
\uncertain{ou encore}~: la donnée d'une structure de $\mathbb{Z}_n$-torseur sur $S$ ($n = \mathrm{card}\, S$). On trouve l'opération $\Psi_r$ ($r \in \mathbb{Z}_n^{*}$)
\[
  (25)\qquad \Psi_r(u) = u^{r}
\]
sur l'ens.\ des structures pol.\ orientées sur $S$, donc on a une opération de $\mathbb{Z}_n^{*}$ sur l'ens.\ $\mathrm{Polor}(S)$ des structures de polygone orienté sur $S$. Ici cette opération est encore libre.

Si $P = (E, S, \omega)$ est un $n$-polygone régulier, on définit
\[
  (26)\qquad \Psi_r(E) = (E, S, \Psi_r(\omega)) \qquad (\text{si } (r, n) = 1)
\]
qui a encore un sens \add{\ill{}} \struck{\ill{}} \uncertain{c'est un polygone} de $E$, de \uncertain{même} \struck{\ill{}} \ill{} que partie de $E$, mais avec une autre structure combinatoire dessus. Ici \struck{$\mathbb{Z}/\pm 1$} \add{\uncertain{$\Psi_n$}} opère librement sur l'ens.
\[
  \mathrm{Polreg}_n(E)
\]
des $n$-polygones réguliers plongés de $E$. \add{De même}, $D = \mathbb{Z}_n^{*}$ opère \add{librement} sur l'ens.\ des $n$-polyg.\ rég.\ orientés. Item pour pol.\ rég.\ \uncertain{épinglés}
\note{en (26) le premier membre est bien écrit $\Psi_r(E)$, pour $\Psi_r(P)$.}

\page{36}
(plutôt \uncertain{pour les} $n$-pol.\ rég.\ orientés munis d'un sommet~!)
\struck{Enfin, condition $\Pi$ d'unimodularité \ill{} invariante \ill{} $n$-polygone} \struck{combinatoire,}

\textbf{NB} Si on a une \struck{repr.} réal.\ géom.\ régulière d'un $n$-polygone comb.\ $\Pi = (S, \omega)$ ($n$ fini), alors $\Psi_r$ sera la repr.\ de $\Psi_r(\Pi) = (S, \Psi_r(\omega))$ définie par la même injection $S \longrightarrow E$ --- la condition de régularité sera en effet invariante par passage de $\omega$ à $\Psi_r(\omega)$, \struck{\ill{}} étant donné que \struck{le groupe des} l'on a évidemment
\[
  (27)\qquad \mathrm{Aut}(\Pi) = \mathrm{Aut}(\Psi_r(\Pi))
\]
en tant que s-groupes de $\mathfrak{S}_S$, i.e.\ les normalisateurs dans $\mathfrak{S}_S$ de $\omega$ \struck{est \ill{}} \struck{$\Psi_r$} \add{$\Psi_r(\omega)$} sont évidemment les mêmes (ce qui normalise $\omega$ normalise $\omega' = \Psi_r(\omega)$, et inversement car on aura $\omega = \Psi_{r'}(\omega')$, si $r'$ est un inverse de $r$ dans $\Psi_n$).

\textbf{Question} Peut-on avoir \emph{isomorphie} entre un $n$-polygone \add{régulier} $P$ et $\Psi_r(P)$, si $r \in \Psi_n$, $r \neq 1$ dans $\Psi_n$ (i.e.\ si $r \in \mathbb{Z}_n^{*}$, $r \neq \pm 1$ dans $\mathbb{Z}_n^{*}$)~? On trouvera

\page{37}
en termes de l'invariant $\alpha$ que \struck{\ill{}} l'isomorphie entre $P$ et $\Psi_r(P)$ équivaut~: une condition de la forme $\Psi_r(\alpha) = \alpha$ [où $\Psi_r \in \mathbb{Z}[\ldots]$ \ill{} \ill{} \ill{} p.l.\ \uncertain{conn.}, \ill{} \uncertain{on dirait} que \ill{}, \uncertain{pour} \ill{} …] On trouvera \struck{\ill{}} que si $(n, p) = 1$ on aura $\Psi_r(\alpha) \neq \alpha$ si $r \neq 1$ dans $\Psi_n$, tandis que si \struck{\uncertain{complexe}} $p \mid n$ \ill{} (\uncertain{ou bien} $n = p$, $p$ impair, ou bien $n = 2p$) alors \struck{tous} les $n$-pol.\ réguliers sont \uncertain{non isom.}

\struck{Mais} \add{Dans} \uncertain{ces} cas où \struck{\ill{}} \add{\uncertain{puisque}}
\[
  (\mathbb{Z}/p\mathbb{Z})^{*} \simeq \mathbb{Z}/(p-1)\mathbb{Z}
\]
\add{i.e.\ $r_n \neq 1$} \uncertain{est d'ordre} $> 2$ \uncertain{sauf si} $p = 2, 3$ \struck{(\ill{} \ill{} \ill{} \ill{} $p = 2$ \ill{} \ill{} \ill{})}, \ill{} (\uncertain{ce qui correspond} à l'hexagone) la réponse est \uncertain{connue} et \uncertain{négative}, sauf dans les \uncertain{seules} \ill{} \struck{\ill{}} \uncertain{trois} cas \uncertain{analogues} \add{\uncertain{à l'hexagone}} (\uncertain{cas} \ill{} \ill{}~?).
\note{passage très rapide, plusieurs fois récrit ; seuls les énoncés mathématiques sont sûrs.}

On peut retenir quand \uncertain{même} que $\Psi_r$ a une nette tendance~: changer la classe d'isomorphie d'un $n$-polygone régulier. \add{\uncertain{voir} \ill{} \ill{}}

On va définir les $\Psi_r$ ($r \in \mathbb{Z}$) \emph{sans restriction sur $r$} pour les repr.

\page{38}
géom.\ des $\infty$-polygones combinatoires épinglés $\Pi_0(\infty)$. Pour ceci, on note que \struck{la donnée de la \ill{} $\Pi_0(n)$}, \struck{\ill{} \ill{} \ill{}, \ill{} \ill{}} \struck{la repr.\ \ill{}} (pour $n$ quelconque, fini ou infini) la donnée de la repr.
\[
  D_n \xrightarrow{\ \varphi\ } \mathrm{Aff}(E)
\]
équivaut aussi~: la donnée de
\[
  (28)\qquad
  \begin{cases}
    \sigma_E,\ u_E \in \mathrm{Aff}(E) \\
    \text{satisfaisant} \\
    \sigma_E^{2} = \mathrm{id}_E ,\quad \sigma_E u_E \sigma_E^{-1} = u_E^{-1}
  \end{cases}
\]
et si $n$ fini $u_E^{n} = \mathrm{id}_E$,
tandis que la donnée $s_0 \in E$ \struck{satisfait} et les conditions $(2^{\circ})$ et $(3^{\circ})$ \struck{\ill{}} de la prop.~5 \add{\uncertain{détaillées par}} \struck{\ill{} \ill{} \ill{}} \struck{par} le degré de $u_E$ en $u_E^{r}$ --- qui a un sens quel que soit $r \in \mathbb{Z}$, et (dans $n$ fini) \uncertain{conserve}~: simplifie la condition $u_E^{n} = \mathrm{id}$ (après remplacement de $u_E$ par $u_E^{r}$). \struck{\ill{}} Quant~: la condition $2^{\circ}$) \uncertain{de non-alignement} de $s_0$, $u^{r} s_0$ et $u^{2r} s_0$ (\uncertain{ou} de tous les $u^{ir} s_0$, $i \in \mathbb{Z}$, ce qui revient au même)
\note{il écrit ici « condition 2°) » pour la condition de non-alignement, numérotée (3°) dans la prop.~5.}

\page{39}
elle peut ne pas être satisfaite, \add{\uncertain{en fait si elle est satisfaite par $u_E^{r}$}} \struck{elle} dépend des choix de $u_E$ et $s_0$ ($\sigma_E$ est ici indifférent~: une droite) --- donc des « invariants » $\alpha$. \struck{\uncertain{On a, }\ill{}} \add{Par ex.\ si $n = \infty$ \uncertain{mais si}} $\alpha$ est tel que la repr.\ de $D_\infty$ se factorise par $D_n$, donc $u^{n} = \mathrm{id}$, cette condition ne sera pas satisfaite.

Néanmoins on a \uncertain{vu}, \add{\uncertain{\ill{}}} que $\Psi_r(\alpha)$ --- il \uncertain{n'a} pas à être défini par \emph{la} $\alpha$ --- il se pose donc la question de la signification géom.\ de \struck{$\Psi_r(\alpha)$} la repr.\ de $D_\infty$ associée à $\Psi_r(\alpha)$, dans le cas « exceptionnel » $\alpha, r$~! [\struck{\ill{}} \add{\ill{} \ill{} \ill{}, \uncertain{pour} $\alpha$ \ill{} \ill{}} où les $s_i = u^{i} s_0$ sont tous alignés …].
\marginal{\ill{} \ill{} \ill{}}
\note{l'insertion entre crochets est écrite en très petit sur deux lignes ; la marge gauche porte quelques mots en oblique, non lus.}

\section{VII. \struck{Polygones} \add{Systèmes} réguliers \struck{\ill{}} \add{de \ill{}}}

\marginal{(i.e.\ structures polygonales modulo étoilage …)}

Soit $S$ un ensemble \add{fini, de cardinal $n$ (combinatoire)}. On appelle \emph{structure quasi-polygonale} sur $S$ la donnée d'une structure polygonale

\page{40}
sur $S$ modulo une des opérations $\Psi_r$ ($r \in \mathbb{Z}_n^{*}$, ou $r \in \Psi_n$)~; \struck{Je dis \ill{}} en d'autres termes, c'est une orbite \add{$\pi$} de $\Psi_n$ opérant sur $\mathrm{Pol}(S)$. Regardant le s-groupe \add{$Z$ de $\mathfrak{S}_S$} engendré par un $\omega \in \pi$ (lequel ne dépend que de $\pi$, \struck{étant} invariant en remplaçant $\omega$ par $\Psi_r(\omega)$), on trouve un s-groupe cyclique d'ordre $n$, et $S$ est un torseur sous ce groupe. Ainsi, \add{la donnée d'}une structure quasi-polygonale sur $S$ équivaut~: la donnée d'une structure de torseur sur $S$~: groupe \add{$Z$} \struck{des} de translations cyclique (d'ordre $n$ nécessairement). Se donner un $\omega \in \pi$ revient au même que de se donner un gén.\ « au signe près » dans $Z$. Se donner une structure polygonale \emph{orientée} compatible avec la structure quasi-polygonale revient~: la donnée d'un gén.\ de $Z$.

La catégorie des \struck{quasi} quasi-polygones combinatoires d'ordre $n$ ($n \in \mathbb{N}$, $n \geq 3$) est équivalente à celle des \uncertain{fibrés} \ill{} \uncertain{affines} sur $\mathrm{Spec}(\mathbb{Z}/n\mathbb{Z})$, à un tel
\note{la phrase se poursuit au-delà de ce lot.}

\end{document}
